% Task 3 — Demand Mapping and Region Clustering
% Northeast US Megaregion Freight Network Design

\section{Region Clustering}

The third task of the network design pipeline establishes the territorial structure of the
Northeast megaregion by (1) constructing a county-level freight demand surface that guides
spatial analysis and (2) partitioning the 434 county-equivalent units into 50 demand-balanced,
geographically contiguous, and infrastructure-aligned freight regions. These regions form the
fundamental geographic unit for all downstream hub location, gateway design, and flow
assignment work.

% ------------------------------------------------------------------
\subsection{Demand Surface Construction}

\subsubsection{Commodity Filtering and Demand Metrics}

The demand surface is derived from the county-to-county FAF freight flow matrix described in
Task~1 (\texttt{raw.parquet}, 33.4M records). Non-palletizable bulk commodities --- gravel and
mining products (\texttt{sctg1014}) and coal and energy products (\texttt{sctg1519}) --- are
excluded from all demand computations, retaining only truck-compatible freight across the
three palletizable commodity groups. This filter reflects the Physical Internet focus on
unit-load capable flows and is applied consistently through all downstream tasks.

Two county-level demand metrics are constructed from the filtered flow table and serve
distinct roles:

\textbf{Activity throughput.} The bidirectional all-flow endpoint sum --- $\text{tons\_in} +
\text{tons\_out}$ --- captures total freight activity at a county, regardless of whether the
paired endpoint lies inside or outside the megaregion. This metric is used exclusively for
choropleth visualization and the composite infrastructure map; it is not used as a clustering
weight or hub-sizing input because it double-counts intra-region flows that never generate
inter-hub demand.

\textbf{External throughput.} The NE-boundary demand metric retains only flows where one
endpoint is inside the Northeast megaregion and the other is outside. Formally, for county
$c$:
\begin{equation}
  T_c = \text{ext\_in}_c + \text{ext\_out}_c
\end{equation}
where $\text{ext\_in}_c$ is the 2025 tonnage arriving at $c$ from non-NE origins and
$\text{ext\_out}_c$ is the tonnage departing $c$ to non-NE destinations (both in thousand
short tons). This quantity serves as the SA clustering demand weight $w_c$ and, after
aggregation to regions, as the hub-facing regional demand $T_r$ consumed by the Task~5
capacity formulation.

The distinction between these two metrics is critical. Activity throughput reflects raw
freight intensity and is appropriate for visual communication. External throughput reflects
the demand that regional hubs must actually process --- flows crossing the megaregion
boundary --- and is therefore the correct signal for territorial partitioning aimed at hub
network design.

\subsubsection{Composite Infrastructure Map}

County activity throughput is joined to the 2023 Census cartographic boundary shapefile and
projected to EPSG:9311 (US National Atlas Equal Area) to produce a log-scale choropleth of
freight intensity across the 434-county study area. Interstate segments from the NTAD North
American Roads dataset (US interstates: \texttt{CLASS=1}, \texttt{COUNTRY=2}) and rail
segments from the TIGER/Line national rail network are clipped to the megaregion extent and
overlaid. The three tiers of Task~2 interface nodes --- global (seaports, cargo airports),
continental (US--Canada border crossings), and national (external domestic hubs) --- are
superimposed as scaled markers. The resulting composite map synthesizes freight demand,
corridor infrastructure, and boundary interfaces in a single reference frame, providing the
spatial intuition that motivates both the target region count and the corridor-alignment
objective embedded in the clustering formulation.

% ------------------------------------------------------------------
\subsection{Region Clustering via Simulated Annealing}

\subsubsection{Problem Formulation}

The clustering task is to partition the 434 NE county-equivalent units into $k = 50$
regions such that the resulting partition simultaneously satisfies three design properties:
demand balance across regions, geographic compactness within regions, and alignment to
major interstate and rail corridors. No standard closed-form algorithm optimizes all three
jointly under a contiguity constraint, motivating a metaheuristic approach.

Simulated Annealing (SA) on the county adjacency graph is adopted as the optimization
engine. SA is well-suited to this problem class because (a) the feasible space --- contiguous,
non-empty partitions of a planar graph --- has no tractable convex relaxation, (b) the
move neighborhood (single county reassignments between adjacent regions) admits $O(1)$
incremental objective evaluation, and (c) the annealing schedule allows escape from locally
optimal but globally poor solutions that pure greedy methods typically produce on
geography-constrained partitioning instances.

\subsubsection{County Graph and Corridor Weights}

The county adjacency graph is built using positive shared border length as the adjacency
criterion; point-only contacts (e.g., four-corner touches) are excluded. Known island
cases are patched with documented synthetic links. Each edge $(i,j)$ in the graph is
augmented with an infrastructure weight:
\begin{equation}
  \omega_{ij} = f\!\left(\text{interstate\_km}_{ij},\, \text{rail\_km}_{ij}\right)
\end{equation}
computed from the total length of interstate and rail segments crossing the shared border.
This weight enters the alignment penalty described below: cutting an edge with high
$\omega_{ij}$ --- placing counties $i$ and $j$ in different regions across a major
corridor --- incurs a larger penalty than cutting a low-infrastructure border.

\subsubsection{Objective Function}

The SA energy function $J$ is a normalized sum of three penalties:
\begin{align}
  J &= \lambda_{\text{align}} \cdot P_{\text{align}}
      + \lambda_{\text{compact}} \cdot P_{\text{compact}}
      + \lambda_{\text{balance}} \cdot P_{\text{balance}}
\end{align}

\textbf{Alignment penalty} $P_{\text{align}}$ measures the fraction of infrastructure
weight that crosses region boundaries (cut edges weighted by $\omega_{ij}$). Minimizing
this term encourages regions to be bounded by, rather than bisected by, major corridors ---
the appropriate design intent for a hub network whose traffic follows interstate and rail
arteries.

\textbf{Compactness penalty} $P_{\text{compact}}$ is the normalized within-region centroid
sum of squared distances (SSE):
\begin{equation}
  P_{\text{compact}} = \sum_{r=1}^{k} \sum_{c \in R_r}
    \bigl\| \mathbf{x}_c - \boldsymbol{\mu}_r \bigr\|^2
\end{equation}
where $\mathbf{x}_c \in \mathbb{R}^2$ is the EPSG:9311 centroid of county $c$ and
$\boldsymbol{\mu}_r$ is the demand-weighted centroid of region $r$. This term directly
penalizes geographically elongated or dispersed regions that, while technically contiguous,
would impose large travel distances on freight moving between hubs.

\textbf{Demand balance penalty} $P_{\text{balance}}$ penalizes deviation from the target
per-region external throughput $W^* = \sum_c T_c / k$:
\begin{equation}
  P_{\text{balance}} = \sum_{r=1}^{k} \left( \frac{T_r - W^*}{W^*} \right)^2
\end{equation}
where $T_r = \sum_{c \in R_r} T_c$ is the total external throughput of region $r$.
Demand balance is essential because the downstream hub location model sizes hub capacity
proportionally to $T_r$; severe imbalance would produce regions whose demand either
saturates hub capacity or renders a hub unnecessary.

\subsubsection{Move Validity and Incremental Evaluation}

A proposed move reassigns a single county $c$ from its current region $r_{\text{from}}$ to
an adjacent region $r_{\text{to}}$. Four hard filters determine move validity: (1) $c$ must
share a border with at least one county already in $r_{\text{to}}$; (2) the donor region
$r_{\text{from}}$ must retain at least one county; (3) $r_{\text{from}} \setminus \{c\}$
must remain connected (verified by BFS on the region subgraph); and (4) a geometry screening
proxy is applied but is not treated as the binding 5.5-hour travel time constraint, which is
enforced at the hub-selection stage.

Incremental $\Delta J$ evaluation computes only the change in $P_{\text{align}}$,
$P_{\text{compact}}$, and $P_{\text{balance}}$ attributable to the two affected regions,
reducing each proposal to $O(|R_r|)$ rather than $O(N)$ cost. All computationally expensive
geometry operations are cached before the optimization loop begins; the SA loop itself
operates exclusively on arrays, adjacency lists, and cached region statistics.

\subsubsection{Initialization and Search Strategy}

A two-phase warm start is used. First, weighted $k$-means is applied to county centroids
(weighted by external throughput) to identify 50 dispersed seed counties. Second, contiguous
region growing is performed over the county graph using distance and demand-overshoot
penalties to expand each seed into a spatially connected region. The resulting initial
partition is materially better than a pure BFS or nearest-centroid start, yielding an
initial objective of $J_0 = 1.6205$ and providing the SA search a meaningful improvement
trajectory. The annealing schedule uses empirical temperature scaling calibrated to initial
$|\Delta J|$ samples, multiple restarts with independent random seeds, resumable
checkpoints, and an early-stopping patience criterion. The best solution across all restarts
is retained.

% ------------------------------------------------------------------
\subsection{Clustering Results and Quality Assessment}

\subsubsection{Quantitative Solution Quality}

The final SA solution achieves the following metrics:
\begin{itemize}
  \item Objective value: $J^* = 0.4910$ (initial $J_0 = 1.6205$; 69.7\% improvement)
  \item Demand balance: coefficient of variation $\mathrm{CV} = 0.291$,
        $\mathrm{nRMSE} = 0.291$
  \item Corridor alignment: cut fraction $= 0.320$ (32.0\% of infrastructure-weighted
        border length crosses region boundaries)
  \item Contiguity: all 50 regions verified connected via BFS on the county adjacency graph
\end{itemize}

The 69.7\% reduction in $J$ from initialization to termination confirms that the SA search
escapes the warm-start local optimum substantially. A CV of 0.291 on external throughput
implies that the interquartile range of regional demand spans roughly half the mean --- a
level of balance that is practically acceptable for a 50-region design covering a 14-state,
geographically heterogeneous megaregion where population density and freight intensity vary
by two orders of magnitude between rural Appalachian counties and the dense NJ/NY/PA
corridor.

\subsubsection{Structural Interpretation}

The final 50-region partition is summarized in \texttt{region\_assignment.csv} (434 rows;
columns: \texttt{fips}, \texttt{county\_name}, \texttt{state}, \texttt{region\_id},
\texttt{external\_throughput\_ktons}, \texttt{activity\_throughput\_ktons}) and
\texttt{region\_metrics.csv} (50 rows; columns: \texttt{region\_id}, \texttt{n\_counties},
\texttt{external\_throughput\_ktons}, \texttt{centroid\_x\_m}, \texttt{centroid\_y\_m},
\texttt{sse\_m2}, \texttt{is\_connected}).

Several structural observations emerge from the solution:

\textbf{Size heterogeneity.} Region county counts range from 1 (three single-county regions
corresponding to highly urbanized, high-demand areas in PA, NY, and NJ) to 33 (a large rural
region spanning ME, NH, and VT). Mean region size is 8.7 counties (standard deviation 7.3).
This heterogeneity is a direct consequence of the demand-balance objective interacting with
the underlying freight distribution: high-intensity urban counties generate enough external
throughput to justify their own region, while low-intensity rural counties must be aggregated
into large geographic units to reach a comparable demand level.

\textbf{Corridor alignment.} The NJ Turnpike / I-95 corridor, I-81 in the Shenandoah
Valley, and the I-90 Massachusetts Turnpike align with region boundaries rather than
bisecting them, consistent with the alignment penalty design intent. Regions tend to be
bounded on one or more sides by a major interstate, placing hub candidate sites in the
interior of their served territory rather than at corridor crossings.

\textbf{Compactness vs.\ demand balance tradeoff.} Nine regions covering Appalachian areas
(WV, western VA, and rural NY/VT) are geographically large despite relatively low external
throughput, with maximum pairwise county-centroid spans exceeding 150 miles. These regions
reveal the fundamental tension in the formulation: enforcing demand balance in low-density
areas forces large geographic extents, which raises the within-region SSE penalty. The
solution accepts larger compactness penalties for these regions in exchange for globally
better demand balance across all 50 regions. These large-span regions are explicitly flagged
for special treatment in Task~6, where the 80-mile cross-area distance guideline is relaxed
for an exempt set of nine geographically extended regions.

\textbf{Post-assignment demand correction.} After the final region assignment is fixed, the
external throughput weights are recomputed excluding flows where both origin and destination
fall within the same region. This correction removes same-final-region flows that were
counted as external demand during optimization (when final region membership was unknown) but
do not cross region boundaries in the final solution. The corrected regional demand values
are stored in \texttt{region\_external\_metrics.csv} and serve as the authoritative
hub-facing demand $T_r$ for the Task~5 capacity formulation.
